\documentclass[11pt,letterpaper]{article}
\usepackage[margin=.7in,top=.78in,bottom=.7in]{geometry}
\usepackage[T1]{fontenc}
\usepackage{lmodern,amsmath,amssymb,array,booktabs,fancyhdr,hyperref}
\hypersetup{hidelinks}
\renewcommand{\familydefault}{\sfdefault}
\pagestyle{fancy}\fancyhf{}
\fancyhead[L]{\small Week 74 / Doubling elevators / Adult guide}
\fancyfoot[L]{\small Bellingham Math Circle / Week 74 / GGT74-FAC-v1}
\fancyfoot[R]{\small\thepage}
\renewcommand{\headrulewidth}{0pt}\renewcommand{\footrulewidth}{0pt}
\setlength{\parindent}{0pt}\setlength{\parskip}{6pt}
\setlength{\headheight}{14pt}
\newcommand{\sect}[1]{\par\vspace{5pt}{\large\bfseries #1}\par}
\newcommand{\prob}[1]{\par\vspace{3pt}\textbf{Problem #1}\par}
\begin{document}

\sect{The mathematics first}
The board is an infinite graph with states $(x,h)$, where $x$ is any integer and
$h$ is a nonnegative integer. A horizontal move changes $x$ by $\pm2^h$;
an up or down move changes only $h$ by $1$. Every move costs one. All questions
start at $(0,0)$ and finish at height $0$; the printed window is not a boundary.

For any ground-to-ground trip using at most $N$ moves and reaching maximum
height $H$, at least $2H$ moves are vertical. Its absolute displacement is at most
\[
 (N-2H)2^H,\qquad 0\le H\le\lfloor N/2\rfloor.
\]
This includes leftward motion, repeated climbs and horizontal moves on several
levels. The farthest possible right endpoint for a budget $N$ is exactly the
maximum of this expression: climb to a chosen H, use all remaining moves to the
right there, and descend. Thus reaching $16$ needs \textbf{8 moves}.

Children may discover that a longer-looking detour can save moves, and that
``I found a short trip'' differs from ``no trip can be shorter.'' For destination
$23$, going to $24$ and stepping left takes $10$ moves, whereas every trip with
no left moves needs at least $11$. Overshooting $15$ via $16$ only ties the
no-left optimum; it does not demonstrate a strict saving.

\textbf{Who this fits.} One Grades 4--5 prototype, requiring signed-coordinate
tracking, adding and subtracting the printed step sizes, and a partner who can check a rule
while recording one move word. Powers notation and general lower bounds are
optional adult-supported reasoning. There is no separate younger packet or
tested age-fit claim. This nonnegative-height toy graph is related to exponential
distortion in Baumslag--Solitar groups; it is not their full Cayley graph.

\sect{Materials and preparation}
Use the older table of three children as a moving pair and rotating referee,
with one anchored adult. Per pair: one four-page student packet, a spare board
page, pencil and eraser, and a very small pointed marker or pencil tip. Keep one
move word as the recoverable record; count its letters after the trip.

Print at actual size. Adjacent coordinate dots are $0.23$ inches, about
$5.84$ mm, apart. A normal broad pawn can obscure its position; check the chosen
marker before the session. Rehearse U/D at an odd coordinate and R/L on an upper
level. For example, from $(1,1)$ an R goes to $(3,1)$; at $(1,2)$ it goes to
$(5,2)$. Vertical movement must preserve $x$ even when its label is not printed.
Use an extra sheet if a child goes beyond the window; do not forbid the move.
No physical marker-fit or partner-operation rehearsal has yet occurred.

\sect{Launch and a feasible first visit}
Let children move the marker briefly. Show the student example beginning at
coordinate $1$, level $0$: U reaches $(1,1)$, R reaches $(3,1)$, D reaches
$(3,0)$. The recorder writes URD, three moves. Reset to $(0,0)$ for the tasks.
Have a child check a second legal attempt before the pair works independently.

Allow about 10 minutes for Problem 1 and 15--20 for Problem 2, exchanging roles
after a trip. Stop the first exploration with an eight-move route to $16$ and
an explicit distinction between ``best found'' and ``proved best.''
Problem 3's budgets can supply a later proof investigation before Problem 4.
Problem 5, especially the no-left lower bound for $23$, is a readiness-dependent
continuation or return visit; completing every page is not the session goal.
\newpage
\sect{Answers and optional hints}
Each word below begins at $(0,0)$ and finishes on the ground. Other shortest
words are welcome; replay them against the move rule rather than matching spelling.

\prob{1}
\textbf{Coordinate 3 takes 3 moves: RRR. Coordinate 7 takes 6: URRRDR.}
Two moves cannot reach $3$: a ground trip stays at level zero unless its entire
budget is spent going up and down, so its displacement is at most $2$.
Five moves reach at most $6$, by the height bound with $H=0,1,2$ giving
$5,6,4$. Thus the proposed trips to $3$ and $7$ are optimal.
\textit{Hint:} ask what the elevator ride costs before counting the longer strides.

\prob{2}
\textbf{Two optimal eight-move words are UURRRRDD and UUURRDDD.}
They use four strides of $4$ and two strides of $8$ respectively.
Problem 4 proves the lower bound. Do not reject different routes just because
they use intermediate heights; a valid shorter claim should be replayed.

\prob{3}
\begin{center}
\begin{tabular}{@{}ccl@{}}\toprule
Budget & Farthest coordinate & One attaining word\\\midrule
4 & 4 & RRRR\\
5 & 6 & URRRD\\
6 & 8 & UURRDD\\
7 & 12 & UURRRDD\\
8 & 16 & UUURRDDD\\\bottomrule
\end{tabular}
\end{center}
For every permitted H, at most $N-2H$ moves remain horizontal, each worth no more
than $2^H$. Maximize $(N-2H)2^H$ over these finitely many heights. The ascent,
right-strides, descent construction attains every corresponding bound, so the
largest bound is the actual optimum. Using fewer than N moves does not evade
the bound. \textit{Hint:} ``If the highest floor is H, which moves must already
have been spent getting there and home?''

\prob{4}
\textbf{No trip of seven or fewer moves can reach 16 and return to ground.}
For seven moves, the complete height cases are
\begin{center}
\begin{tabular}{@{}ccccc@{}}\toprule
Maximum height H & 0 & 1 & 2 & 3\\
Displacement bound $(7-2H)2^H$ & 7 & 10 & 12 & 8\\\bottomrule
\end{tabular}
\end{center}
All are below $16$; reaching height $4$ and returning already costs $8$ moves.
The triangle inequality is the key to the generality: the absolute value of a
sum of signed strides is at most the sum of their absolute values. Repeated
climbs use \emph{more} vertical moves, intermediate floors give \emph{smaller}
strides, and left steps cannot exceed the absolute-value allowance. None makes
a loophole. Together with Problem 2, this proves an exact optimum of $8$.

\sect{Keep the argument concrete}
Before writing powers, children can set aside two move tokens for every level
climbed, then ask how far each remaining horizontal token could possibly go.
These tokens represent an upper bound, not a requirement that all actual trips
use only one level. Let the partner challenge a bound with a route and explain
which allowance it uses. This preserves children's route choices while making
the exhaustive height cases visible.
\newpage
\prob{5}
\begin{center}
\begin{tabular}{@{}cclc@{}}\toprule
Target & Minimum & One shortest word & Minimum with no L\\\midrule
9 & 7 & UURRDDR & 7\\
15 & 9 & UUURRDDDL & 9\\
17 & 9 & UUURRDDDR & 9\\
23 & 10 & UUURRRDDDL & 11\\\bottomrule
\end{tabular}
\end{center}
\textbf{Lower bounds for 9 and 17.} Six moves reach at most $8$, and eight
moves reach at most $16$ (Problem 3), so these targets need at least $7$ and $9$.

\textbf{Why 15 needs 9.} With budget $8$, heights $0$ and $1$ allow only $8$
and $12$ displacement; height $4$ spends all moves vertically. At height $2$,
at most four horizontal moves remain. An odd endpoint needs an odd stride,
which is possible only on level $0$, leaving displacement at most
$1+4+4+4=13$. At height $3$, only two horizontal moves remain; the analogous
bound is $1+8=9$. These are all cases, so eight moves cannot reach $15$.

\textbf{Why 23 needs 10.} For budget $9$, the height bounds for
$H=0,1,2,3,4$ are $9,14,20,24,16$. Only $H=3$ could reach $23$.
It permits at most three horizontal moves; one must be a ground stride to
produce an odd endpoint. Their absolute displacement is at most $1+8+8=17$.
So nine moves cannot suffice. The displayed ten-move word reaches $24$ at
height zero and then steps left to $23$.

\textbf{The no-left comparison needs its own proof.} If the largest visited
height is H, rightward strides are coins of values $1,2,\ldots,2^H$.
The fewest coins totaling an integer target $n$ is
\[
 c_H(n)=\lfloor n/2^H\rfloor+\operatorname{ones}(n\bmod2^H),
\]
where $\operatorname{ones}$ counts the $1$s in a binary expansion. Combining
two equal smaller coins into the next size shows that a minimum collection has
at most one of each size below $2^H$; the remaining smaller coins are precisely
the binary remainder. Every trip therefore needs at least $2H+c_H(n)$ moves.
This bound is attained by ascending to H and taking the required strides while
descending through their levels.

For $n=23$, the bounds at heights $0,1,2,3,4$ are
$23,14,11,11,12$. Heights at least $5$ already use ten vertical moves and need
at least one horizontal move, so cannot beat $11$. The eleven-move no-left
word UUURRDRDRDR gives $16+4+2+1=23$; the ten-move word with L is strictly better.
For the other targets, no-left witnesses are UURRDDR for $9$,
UURRRDRDR for $15$, and UUURRDDDR for $17$. Since these attain the unrestricted
minima, leftward motion cannot improve them. \textit{Hint:} ``Can passing the
target by one help? Does it improve the best no-left route, or merely tie it?''

\sect{Source lineage and status}
Nicholas Touikan, \emph{An Introduction to Combinatorial and Geometric Group
Theory}, \href{https://ntouikan.ext.unb.ca/MATH6022/IntroCGGT/html_output/section-15.html}{Section 3.2.3, HNN extensions and dual tracks}, discusses
$BS(1,2)=\langle a,t\mid t^{-1}at=a^2\rangle$ and its distorted cyclic subgroup.
This supplies the exponential-shortcut connection. The stated graph, route
instances and optimality proofs here are independently developed, not a literal
Cayley-graph activity taken from that source.

\textbf{Review status:} matched to the final four-page GGT74-S-v1 packet.
An independent, coordinate-unbounded finite-depth search checks every numerical
answer and route; the proofs above establish the all-routes conclusions.
The guide is built, rendered and visually checked. It remains an unpiloted
prototype with physical setup and classroom pacing unrehearsed.

\end{document}
